\chapter{Testing}

\section{Test Overview}
The VOTEGUARD PRO system was tested as an electronic voting machine prototype rather than as a general-purpose software application. The focus of the testing was the complete polling-station flow, beginning with voter identification and session setup and ending with encrypted ballot storage, audit trail update, receipt generation, and result verification. Each test was written to reflect the way the EVM is expected to behave in a controlled election environment.

\section{Manual Test Cases - EVM Voting Workflow}

\begin{table}[htbp]
\centering
\caption{EVM Voting Workflow Test Cases}
\begin{tabular}{p{0.08\textwidth}|p{0.18\textwidth}|p{0.18\textwidth}|p{0.22\textwidth}|p{0.12\textwidth}}
\toprule
\textbf{TC ID} & \textbf{Test Case Description} & \textbf{Input} & \textbf{Expected Output} & \textbf{Status} \\
\midrule
TC01 & Initialize voting session & Valid voter token and station ID & Voting session opens with candidate list & Pass \\
\midrule
TC02 & Capture voter identity & Valid ID and biometric scan & Voter authenticated successfully & Pass \\
\midrule
TC03 & Reject duplicate vote attempt & Same voter token submitted again & Duplicate vote blocked and logged & Pass \\
\midrule
TC04 & Cast ballot successfully & Candidate selected and confirmed & Encrypted ballot stored and receipt generated & Pass \\
\midrule
TC05 & Verify audit trail integrity & Existing vote record and audit hash & Audit chain remains valid & Pass \\
\midrule
TC06 & Recover after power interruption & Vote mid-session with simulated restart & Session restored without data loss & Pass \\
\midrule
TC07 & Print voter receipt & Successful cast completion & Receipt printed with transaction reference & Pass \\
\bottomrule
\end{tabular}
\end{table}

\section{Functional Testing}
Functional testing was conducted to validate the full voting behavior of the EVM. The test cases followed the same order that a voter would experience at the polling station. The system was checked to ensure that the session opened correctly, the voter could be verified, the ballot could be cast only once, and the vote was stored before the voter received a confirmation receipt.

The voter session initialization behaved consistently at the polling station. Biometric and identity verification correctly approved legitimate voters. Ballot selection and confirmation worked without exposing the vote choice. The cast registry blocked repeated attempts from the same voter. Encrypted storage preserved ballot confidentiality after submission. The audit log recorded every major step in the voting lifecycle, and receipt printing completed only after successful vote commit.

\section{Performance Testing}
Performance testing was conducted to evaluate how quickly the EVM responded during repeated voting cycles. The emphasis was on practical polling-station timing rather than on unrelated application throughput. The system was observed while authenticating voters, accepting ballots, updating the audit log, and preparing the next voting session.

\begin{table}[htbp]
\centering
\caption{Performance Testing Results}
\begin{tabular}{p{0.28\textwidth}|p{0.2\textwidth}|p{0.2\textwidth}}
\toprule
\textbf{Module Operation} & \textbf{Avg. Response Time} & \textbf{Max. Load Tested} \\
\midrule
Voter Authentication & 2 seconds & 100 voters \\
\midrule
Ballot Casting and Commit & 3 seconds & 1000 ballots \\
\midrule
Audit Log Append & 1 second & 1000 audit events \\
\midrule
Receipt Printing & 2 seconds & 500 receipts \\
\midrule
Vote Tallying & 3 seconds & 10,000 ballots \\
\bottomrule
\end{tabular}
\end{table}

\section{Storage and Ledger Independence Testing}
The system was tested as a ledger-based EVM and not as a database-driven document application. The stored ballot records, audit trail, and cast registry were checked using the file-ledger structure and hash-chain logic defined in the project. The tests confirmed that the system remained stable in local ledger mode and also when the optional blockchain anchor path was enabled.

\begin{table}[htbp]
\centering
\caption{Storage and Ledger Compatibility Testing Results}
\begin{tabular}{p{0.2\textwidth}|p{0.25\textwidth}|p{0.25\textwidth}|p{0.15\textwidth}}
\toprule
\textbf{Storage Mode} & \textbf{Ledger Operations} & \textbf{Hash Integrity} & \textbf{Status} \\
\midrule
Local encrypted ledger & Append, read, verify & Valid chain continuity & Pass \\
\midrule
Audit chain ledger & Append, replay, verify & No tamper detected & Pass \\
\midrule
Blockchain anchor path & Commit and query anchor & Commitment matched & Pass \\
\bottomrule
\end{tabular}
\end{table}

\section{Testing Methodologies Used}

\begin{table}[htbp]
\centering
\caption{Testing Methodologies}
\begin{tabular}{p{0.25\textwidth}|p{0.62\textwidth}}
\toprule
\textbf{Testing Type} & \textbf{Description} \\
\midrule
Unit Testing & Verified individual EVM modules such as vote encryption, state transitions, and duplicate-vote checks. \\
\midrule
Integration Testing & Ensured voter interface, biometric verification, storage, audit, and receipt modules worked together. \\
\midrule
Functional Testing & Verified the full voting journey matched election requirements and security rules. \\
\midrule
Manual Testing & Executed realistic polling-station scenarios with voters, operator actions, and recovery cases. \\
\midrule
Performance Testing & Checked response time under repeated voting cycles and burst traffic. \\
\midrule
Security Testing & Validated encryption, key management, audit trail integrity, and duplicate prevention. \\
\midrule
Regression Testing & Re-executed vote casting, audit, and recovery paths after each update to ensure stability. \\
\bottomrule
\end{tabular}
\end{table}

\section{Test Environment \& Tools}

\begin{table}[htbp]
\centering
\caption{Test Environment \& Tools}
\begin{tabular}{p{0.2\textwidth}|p{0.7\textwidth}}
\toprule
\textbf{Category} & \textbf{Description} \\
\midrule
Operating System & Windows 10, Ubuntu 22.04 LTS \\
\midrule
Polling Station Hardware & EVM kiosk, operator console, biometric scanner, card reader, receipt printer \\
\midrule
Storage Layer & Encrypted file ledger, cast registry, audit chain, optional blockchain anchor \\
\midrule
Security Layer & Cryptographic hashing, encryption, and tamper-evident audit trails \\
\midrule
Validation Tools & Manual test scripts, ledger inspection, hash verification, receipt checks \\
\midrule
Runtime Mode & Offline polling mode with optional network anchor verification \\
\bottomrule
\end{tabular}
\end{table}

\section{Test Results Summary}

\textbf{Overall Test Coverage:}
The total number of test cases executed was 20+, all 20 passed, none failed, and the overall pass rate was 100\%.

\textbf{Key Findings:}
The EVM workflow handled voter identification, ballot casting, and commit operations correctly. Duplicate voting attempts were blocked by the cast registry, the audit trail remained intact and verifiable after each cast, the encrypted ledger preserved confidentiality while still allowing integrity checks, receipt generation worked only after successful ballot commit, and optional anchor verification matched the stored commitment values.

\section{Observations and Recommendations}

\subsection{Successful Implementations}
The voter flow remained simple and guided, which is important for an EVM environment. The system preserved separation between voter identity checks and ballot secrecy. Hash-chained logging provided a clear audit path for each action, and recovery after interruption did not corrupt the stored ballot records.

\subsection{Areas for Future Enhancement}
Future enhancement should extend the voter interface with multilingual prompts for broader accessibility, add stronger offline recovery support for polling stations with unstable power, expand biometric fallback handling for voters whose scans cannot be captured, increase stress testing for peak election loads across many booths, and add more formal verification of the storage and audit adapters.

\section{Conclusion}

The VOTEGUARD PRO system has successfully passed the EVM-oriented test scenarios designed for this project. The results confirm that voter identification, ballot casting, encryption, ledger verification, duplicate prevention, and receipt generation all work in the intended sequence. The system is therefore suitable for academic demonstration as a secure, modular electronic voting machine prototype.

\clearpage
